%======================================================================
\subsection{PDM-14 and SIM-07 follow-ups --- Built, 1 Nov}\label{sec:followups-built}
%======================================================================

Seven follow-ups you approved after the PDM-14 and SIM-07 phase~1 reports, in two commits (PDM-14: four items;
SIM-07: the named-selection and Mesher items). The spin-then-vibrate validation case was \emph{not} built: see
the open questions.

\paragraph{What you can do now.}
\begin{itemize}[nosep]
  \item \textbf{Save As lists the simulation filter.} When the active document is a simulation (its file ends in
        \code{.vlsim}, or it holds the \code{Simulation source} object of a study copied from a part), the Save As
        dialog lists ``Vector Labs simulation (*.vlsim)'' first and FCStd second; a normal document still sees
        FCStd only (\file{src/Gui/Document.cpp}, \code{Document::saveAs}). A name typed without a suffix takes
        \code{.vlsim} from the first filter. The native dialog cannot be driven by a test: the C++ builds, the
        filter logic is two lines.
  \item \textbf{Refresh Pictures records the shape fingerprint} of a part, assembly or bought item that has none, in
        the same Mac-side cache as the pictures (never in the vault), so studies of parts checked in before PDM-14
        stop reading ``unknown''. A part with no solid records an empty fingerprint, so it is not opened again.
        An assembly drawn from an older work copy is skipped (that file is not this version's shape).
  \item \textbf{Parts with several bodies:} the fingerprint covers all the visible solids (a compound: volumes and
        areas summed, the box of all, face and solid counts summed). One body gives exactly the shape it always gave;
        a test pins the box value recorded before the change (\code{7c4ec29572c5}).
  \item \textbf{All projects} no longer shows the four totals tiles when project cards are shown. A vault without
        projects keeps them. The tiles were in \file{bruceworx\_dashboard.py}, not in the project page (whose own
        count tiles stay).
  \item \textbf{A load, support or local size may use several named selections plus faces of its own.} Its faces
        are the union, duplicates removed. Settings: \emph{Named selections} (comma separated) and \emph{Extra
        faces}; the single \emph{Named selection} pick still works and reads ``(several)'' when there are more.
        Editing a selection moves every user and the solution goes stale, as before; removing one of several keeps
        its faces as the object's own. Older documents (one \code{BW\_Selection} link) read as a list of one.
  \item \textbf{Mesher} (Netgen or Gmsh) is a setting of the Mesh node beside the time limit; Netgen is the
        default and leaves every existing fingerprint unchanged; Gmsh marks the mesh and the solution stale.
\end{itemize}

\paragraph{Decided while building.}
\begin{itemize}[nosep]
  \item The ``refresh'' fingerprint is cached by the version's hash, like the summary, so it costs one file open.
  \item \code{Mesher} is left out of the study fingerprint and put into the mesh fingerprint only when it is not
        Netgen: adding the property to an older study does not make it stale.
  \item Extra faces cannot be set on an object that uses no named selection (it has its own list already).
\end{itemize}

\paragraph{Tests.} Addon unit tests 730 OK; headless TestBruceWorxSimModel 50 OK, Preview 12 OK, Vlsim 2 OK; GUI TestBruceWorxSimWorkspace 32 OK, TestBruceWorxProjectPage 9 OK. Besides that, \code{TestBruceWorxSimVault} has one error that
predates this work (\code{test\_results\_summary\_and\_stripped\_copy}, GUI run: the stripped copy loses
\code{max\_vm} after a re-save; it fails the same with the earlier module).

\paragraph{Open questions.}
\begin{itemize}[nosep]
  \item \textbf{Spin, then vibration modes.} I found no closed form for a spinning disc I could cite with
        confidence: Lamb and Southwell (1921) and Southwell (1922, Proc.\ R.\ Soc.\ A 101, 133--153) are the
        sources, but their formulas are not available to me as text. Not guessed. Options: you supply the paper, or
        I use the rotating cantilever beam (Wright et al., 1982) as the benchmark instead.
\end{itemize}
